\documentclass{scrreprt}

\usepackage[utf8]{inputenc}
\usepackage[francais]{babel}
\usepackage[T1]{fontenc}
\usepackage{lmodern}

\usepackage{graphicx}

\begin{document}

\thispagestyle{empty}

\phantom{~} % lolilol

\vspace{7em}

\hrule

\vspace{0.5em}

\begin{center}
	\Huge \textbf{Rapport technique} \\
	\LARGE Projet \og IUTBM\fg\ 
\end{center}

\vspace{0.5em}

\hrule

\vspace{5em}

\begin{center}
	\large
	\begin{tabular}{rl}
	 Marine & \bsc{Benoît} \\
	Florian & \bsc{Durosier} \\
	  David & \bsc{Nowinski} \\
	    Ali & \bsc{Sid} \\
	 Julien & \bsc{Voisin} \\
	 Cédric &\bsc{Zambelli}
	\end{tabular}

	\vspace{4em}

	\normalsize
	Vendredi 23 mars 2012 \\

	\vspace{2em}

	\includegraphics[width=2cm]{logoiut.eps}

	Institut Universitaire de Technologie de Belfort-Montbéliard \\
	Département Informatique
\end{center}

\newpage

\section*{Introduction}

Pour réaliser le projet IUTBM, nous avons choisi le langage Python, qui permet une
 programmation plus aisée des différents algorithmes. Nous avons choisi le module
 Pygame (basé sur la bibliothèque SDL) pour la réalisation de l'interface graphique.


\end{document}
